\chapter{Conclusion and Future Work}
\label{cha:Conclusion}


todo sollten eigene Absätze sein, statt unterkapitel
\section{Future Work}

Whilst the work done in this thesis serves as a good starting point, it is possible to branch off in many different directions for further exploration of the usages for real-time inpainting
and clutter removal. Whilst the approach for area selection, as presented in \ref{cha:Concept} is easy to use and serves its purpose, there are several advancements to automatically detect clutter based
on a users selection. In 2023 an article by Huynh, Zhou et. al. \cite{HuynhChuong2023SSPC} proposed a methodology to easily detect similar object patterns within an image. This may be used to allow for the user to simply select an object by pointing at it, followed by the SimpSON algorithm creating masks around the image for the necessary inpainting.

In addition to this, it might be interesting to conduct a user study on the effects of hiding distractions using inpainting in AR, how it affects the user and if there are any differences to other methods of hiding objects. Whilst the work of Cheng, et. al \cite{ChengYiFei2022TUDR} has already laid good groundwork on the effects of DR, inpainting was not considered at that time
and may outperform the approaches which were analyzed.

The proposed and implemented evaluation prototype was built in an easily extendable manner, therefore it would be interesting to evaluate more methodologies and algorithms, as there are still many promising algorithms which could not be analyzed in further detail, not to mention newer algorithms and methodologies which might be published over the next years.

As this thesis was heavily focussed on CPU based methodologies, converting these algorithms into GPU-compatible versions and comparing these should definitely be addressed in the future to ensure the best performance of the algorithms.